\documentclass[11pt]{article}

% Change "review" to "final" to generate the final (sometimes called camera-ready) version.
% Change to "preprint" to generate a non-anonymous version with page numbers.
\usepackage[preprint]{acl}

% --- IEEE-style numeric citations (override acl.sty's author-year natbib defaults) ---
\setcitestyle{numbers,square,sort&compress}

% Standard package includes
\usepackage{times}
\usepackage{latexsym}

% For proper rendering and hyphenation of words containing Latin characters (including in bib files)
\usepackage[T1]{fontenc}
% For Vietnamese characters
% \usepackage[T5]{fontenc}
% See https://www.latex-project.org/help/documentation/encguide.pdf for other character sets

% This assumes your files are encoded as UTF8
\usepackage[utf8]{inputenc}

% This is not strictly necessary, and may be commented out,
% but it will improve the layout of the manuscript,
% and will typically save some space.
\usepackage{microtype}

% This is also not strictly necessary, and may be commented out.
% However, it will improve the aesthetics of text in
% the typewriter font.
% \usepackage{inconsolata}  % cosmetic typewriter font; requires texlive-fonts-extra

%Including images in your LaTeX document requires adding
%additional package(s)
\usepackage{graphicx}

\usepackage{booktabs}
\usepackage{longtable}
\usepackage{array}
\usepackage{xcolor}
\usepackage{colortbl}
\usepackage{caption}
\usepackage{adjustbox}   % for \begin{adjustbox}
\usepackage{natbib}      % for \citeyear  (or use biblatex)
\usepackage{lscape}      % optional: landscape
\usepackage{multirow}
\usepackage{amssymb}

\title{Methodology on the Multi-Page Visually Rich Document Understanding Frontier: A Literature Review}

\author{Author: Lewei Xu \\
  The University of Western Australia \\
  \textbf{Supervisors}: Dr. Yihao Ding, Dr. Daochang Liu \\
  \textbf{Degree}: BACS (Honours) \\}

\begin{document}
\maketitle
\begin{abstract}
Multi-page visually rich document understanding (MP-VRDU) has emerged as a distinct research area driven by the prevalence of long-form documents such as contracts, scientific papers, and financial reports whose evidence is distributed across pages and modalities. While single-page document understanding has matured into a well-defined sub-field with established benchmarks, the multi-page setting introduces challenges that simple extension cannot address: total document length frequently exceeds backbone context windows, answer-bearing evidence is sparsely distributed, and cross-page structural relationships cannot be recovered from per-token layout signals alone. This review synthesises recent work in MP-VRDU, organising the field by how systems manage cross-page evidence rather than by OCR dependency. We examine four architectural families, three modality-specific challenges, and the inference-time and training strategies that recur across them, before identifying three persistent gaps that current approaches have not resolved: data scarcity coupled with evaluation integrity concerns, the evidence pipeline bottleneck between coverage and fidelity, and cross-page structural reasoning that remains brittle across all current routes.
\end{abstract}

\section{Introduction}
\label{sec:introduction}
Documents are the principal medium through which organisations, governments, and the scientific community record and exchange information. Contracts, financial reports, scientific papers, technical manuals, regulatory filings, and presentation decks are pervasive and structurally heterogeneous, interleaving dense text with tables, figures, diagrams, and complex layouts~\cite{tito2023hierarchical,borchmann2025arctic}. Manually extracting information from such documents is slow, expensive, and error-prone, motivating a sustained research effort to build automated document visual question answering systems that can answer natural-language questions over rendered document pages.

For single-page inputs, this problem is largely solved. Document-specialised systems such as Qwen2.5-VL~\cite{bai2025qwen2} and mPLUG-DocOwl1.5~\cite{hu2024mplug1.5} approach or exceed human performance on DocVQA~\cite{mathew2021docvqa}, InfographicVQA~\cite{mathew2022infographicvqa}, and ChartQA~\cite{masry2022chartqa} datasets. Recent frontier multimodal LLMs such as Claude Opus 4.7~\cite{anthropic2026opus47} and Gemini 3.1 Pro~\cite{google2026gemini31} can now read a rendered page directly from pixels, treating the document as an image and matching or exceeding human reading on single-page and even short multi-page documents provided the input resolution is sufficient. Real-world documents, however, almost always extend beyond a single page, and multi-page visually rich document understanding (MP-VRDU) is not a simple extension of its single-page counterpart.

Real-world documents, however, almost always extend beyond a single page, and the evidence required for a realistic question is typically scattered across several pages or distributed across heterogeneous elements~\cite{tito2023hierarchical,blau2024gram,borchmann2025arctic}. Three distinct challenges emerge when moving from one page to many. First, input length scales with page count, with each visually rich page contributing high-resolution figures, tables, and charts that quickly exceed the context window of existing language and vision-language models; quadratic-cost self-attention~\cite{vaswani2017attention} compounds the compute burden further. Second, the evidence required to answer a question is typically sparse and distributed across pages, requiring a system to navigate, retrieve, and reason across the entire document rather than read a single page in detail \cite{blau2024gram}. Third, explainability becomes more demanding, as models must not only produce an answer but also localise the pages or passages from which it was derived. These pressures are compounded by the well-documented tendency of long-context models to attend poorly to evidence buried in the middle of long inputs \cite{liu2024lost}.

The past two to three years have seen rapid growth in approaches addressing these challenges, with the field broadly evolving from joint multi-page encoding toward increasingly selective and iterative document-level systems. Early work extended single-page document transformers to multi-page inputs through hierarchical or global-reasoning architectures \cite{tito2023hierarchical,blau2024gram}. Later approaches moved toward backbone-centric adaptation of pretrained multimodal large language models (MLLMs) tuned for multi-page reasoning \cite{duan2025docopilot,hu2025mplugdocowl2}. As document length remained a persistent constraint, retrieval-augmented pipelines became more prominent \cite{cho2024m3docrag,tanaka2025vdocrag}, and more recent work has adopted agent-based systems that iteratively acquire evidence rather than retrieving once \cite{sun2025docagent,wang2025vidorag}. This rapid succession has produced inconsistent terminology and limited cross-family comparison, making it difficult to identify which design choices most influence performance in the multi-page setting. Prior surveys have charted document understanding from complementary angles \cite{subramani2020survey,ding2026deep,gbada2025deep,ding2025survey,gao2025scaling}, but remain anchored to the single-page setting or organise progress by OCR dependency, treating multi-page documents as an incremental extension rather than a distinct problem.

This review synthesises recent MP-VRDU research with the aim of identifying persistent gaps that motivate continued research in this area. Section~\ref{sec:background} briefly summarises the single-page mechanisms multi-page systems inherit, orienting the discussion toward the design decisions that the multi-page setting reformulates. Section~\ref{sec:architecture} introduces four architectural families that organise the field through the lens of cross-page evidence management. Sections~\ref{sec:modality} through~\ref{sec:training} narrow into specific design dimensions: multimodal representation under context constraints, inference-time composition strategies, and component-level training. Section~\ref{sec:datasets} catalogues the resources that support training and evaluation, and Section~\ref{sec:gaps} identifies three persistent open gaps in the field.

\section{Background: Single-Page VRDU}
\label{sec:background}

Multi-page VRDU systems inherit most of their per-page mechanisms from the single-page literature, recently surveyed by Ding et al. \cite{ding2025survey}. We summarise the two organising axes that recur into the multi-page setting: the OCR-dependent versus OCR-free split at the architectural level (\S\ref{sec:bg-architecture}), and three orthogonal modality-handling choices that persist independently of that split (\S\ref{sec:bg-modality}).

\subsection{Architectural Families}
\label{sec:bg-architecture}

The dominant axis at the single-page level distinguishes two architectural families that trade different costs. \textbf{OCR-dependent frameworks} preprocess the page with off-the-shelf optical character recognition (OCR) or PDF parsers, producing recognised text and bounding-box layout that is either fed into the language model prompt directly \cite{wang2024docllm,he2023icld3ie} or projected through an auxiliary multimodal encoder such as LayoutLMv3 \cite{huang2022layoutlmv3,luo2024layoutllm}. These systems rely on external tools to capture structural information without extensive visual pretraining, but reliance on OCR introduces cumulative errors, particularly on handwritten or low-quality scanned documents, and the use of low-resolution document images limits the expressiveness of visual representations. \textbf{OCR-free frameworks} discard the parser and have a vision encoder read text directly from rendered page pixels \cite{kim2021donut,ye2023ureader,hu2024mplugdocowl}, with the language model trained to recover text and layout through pretraining objectives such as text recognition, grounding, and captioning \cite{wang2023towards,feng2023unidoc,liu2024textmonkey}. These frameworks enable end-to-end training and handle handwriting and complex layouts that OCR cannot, but accurate fine-grained text recognition demands high-resolution input and large-scale pretraining, which the visual compression and shape-adaptive cropping strategies discussed in \S\ref{sec:bg-modality} are designed to manage. The two families thus trade different costs: OCR-dependent systems inherit cumulative parser errors but avoid expensive visual pretraining, whereas OCR-free systems handle handwriting and complex layouts end-to-end but require high-resolution input and large-scale pretraining.

\subsection{Multimodal Representation}
\label{sec:bg-modality}

Across both architectural families, three orthogonal design decisions govern how text, layout, and visual modalities are encoded, with a fourth cross-cutting choice for how they are fused. \textbf{Text} is recovered as OCR strings injected into the prompt \cite{wang2024docllm,kim2024dockd}, through OCR-augmented multimodal encoders that produce joint text-layout-vision embeddings \cite{luo2024layoutllm,zhu2025enhancing}, or as a training objective for the vision encoder itself through text recognition and detection tasks \cite{wang2023towards,feng2023unidoc,wang2025marten}. \textbf{Layout} is encoded as 2D positional features attached to text tokens \cite{xu2020layoutlm,zhu2025enhancing}, as verbalised bounding boxes or markdown-style formatting in the prompt \cite{he2023icld3ie,lamott2024lapdoc,perot2024lmdx}, as specialised layout tokens that share positional information with their associated text \cite{zhu2025simple}, or through supervised pretraining tasks such as visual text grounding and table reconstruction \cite{liu2024hrvda,liu2024textmonkey,liao2024doclayllm}. \textbf{Visual tokens} are managed through low-resolution patch embeddings \cite{xie2024pdfwukong,tanaka2024instructdoc}, high-resolution shape-adaptive cropping that divides oversized pages into fixed-size sub-images \cite{ye2023ureader,hu2024mplugdocowl}, or compression via resamplers, Q-Formers, and convolutional reducers that reduce the visual token count by an order of magnitude \cite{hu2025mplugdocowl2,yu2024texthawk2,liu2024textmonkey}. \textbf{Modality fusion} proceeds through dedicated neural encoders such as LayoutLMv3 that produce joint multimodal embeddings \cite{huang2022layoutlmv3,luo2024layoutllm}, task-oriented objectives that bind modalities through shared supervision \cite{hu2024mplugdocowl,wang2025marten}, or prompt-based combinations that interleave layout descriptions with text and images at the language model input \cite{lamott2024lapdoc,perot2024lmdx,luo2024layoutllm}. Single-page systems make these choices assuming the entire document fits in a single backbone input, and modality fusion operates on full page-level evidence at once. The multi-page setting reformulates each into a question of which page or region to attend to next, and modality fusion becomes a question of which evidence to combine after selective inspection rather than which to combine within a fixed input.

\section{Architectural Landscape}
\label{sec:architecture}

The single-page literature is naturally organised by OCR dependency, but in the multi-page setting the defining challenge is how a system manages document-scale evidence: which information to retain, which to discard, and how the surviving information should interact across pages. Through this evidence-management lens, the field has shifted from encoding all pages jointly toward increasingly selective and iterative document-level handling, producing four recurring trends: page-to-document encoding models, backbone-centric adaptation, retriever-generator pipelines, and agentic pipelines.

\subsection{Page-to-Document Encoding Models}
\label{sec:arch-p2d}

Page-to-document encoding models process each page through a shared document encoder that fuses OCR text, layout, and visual features at the page level, then connect per-page representations through a document-level exchange mechanism within a single forward pass. The mechanism may take the form of per-page summary tokens \cite{tito2023hierarchical}, document-global tokens with local-global attention \cite{blau2024gram}, recurrent memory states passed across pages \cite{dong2024multipage}, or sparse cross-page attention over long OCR sequences \cite{borchmann2025arctictilt}. By preserving page identity, local layout, and reading order in the architecture itself, this design provides a strong inductive bias for multi-page structure and is well suited to documents with regular structural templates such as multi-page forms and business reports \cite{tito2023hierarchical, borchmann2025arctictilt}. The principal limitation is intrinsic to the design: every page must still be encoded jointly in a single forward pass, computation grows with document length, and the page-to-document compression step inevitably discards fine-grained evidence as inputs scale \cite{blau2024gram, dong2024multipage}. As the context windows of pretrained MLLMs have widened, the practical motivation for a dedicated cross-page encoder has weakened, and subsequent families absorb cross-page information into the backbone or select evidence before reading.

\subsection{Backbone-Centric Adaptation}
\label{sec:arch-bc}

Backbone-centric adaptation removes the dedicated cross-page encoder entirely and processes the document inside a pretrained MLLM whose context window absorbs cross-page information directly. Adaptation strategies span a spectrum of training intensity. At the lightweight end, projectors map per-modality features into the backbone's token space \cite{tanaka2024instructdoc}, adapters enable parameter-efficient fine-tuning \cite{zhu2025simple}, and visual-token compression modules constrain token growth on long inputs to keep multi-page documents within the context budget \cite{hu2025mplugdocowl2}. At the heavier end, the backbone itself is unfrozen for continued pretraining or instruction tuning on document corpora, trading parameter efficiency for stronger document-specific grounding \cite{duan2025docopilot}, with long-context tuning further improving moderate-length multi-page performance \cite{hu2024mplugdocowl, yu2024texthawk2}. The family inherits the foundation model's instruction-following and reasoning, avoids the engineering overhead of external retrieval or orchestration, and remains parameter-efficient to train. The cost of this design is that performance degrades beyond moderate-length documents, leaving quality bound to context length and to how faithfully fine-grained text and visual signals survive token compression \cite{jia2024leopard}. Reliable text recovery depends on high-resolution input that aggressive downsampling readily compromises, and document structure such as section hierarchy and figure-caption linkage must be recovered implicitly from pixels rather than supplied through layout tokens or a structural encoder \cite{hannan2025docslm, duan2025docopilot}. Attention dilution within long contexts \cite{liu2024lost} further weakens evidence buried in the middle, motivating subsequent families that move evidence selection out of the backbone altogether.

\subsection{Retriever-Generator Pipelines}
\label{sec:arch-rg}

Retriever-generator pipelines decouple evidence selection from answer generation. A retrieval module returns a compact subset of OCR chunks, pages, visual patches, multimodal elements, or graph nodes, and a generator answers conditioned only on those units \cite{cho2024m3docrag, tanaka2025vdocrag}. The design adapts the retrieval-augmented generation paradigm from text \cite{lewis2020retrievalaugmented, guu2020realm} to visually rich inputs, with the retriever instantiated as a sparse lexical scorer, a dense embedding model, or a structure-aware index, and the generator typically a pretrained MLLM. Inference cost depends on the retrieved content rather than document length, allowing the architecture to scale to documents that exceed any current MLLM context window. Modular decomposition lets retriever and generator be swapped independently, the retrieved evidence provides a partial audit trail for the answer, and the same interface extends to multi-document corpora without architectural change. Structured domains with reliable retrieval anchors such as scientific papers and technical manuals benefit most \cite{shin2025multidocfusion}, and page-level visual retrieval preserves figures, tables, and layout that OCR-only pipelines lose \cite{tanaka2025vdocrag, xie2024wukong}. The principal limitation is the single-pass control flow, which assumes one retrieval step can surface all evidence required for the answer. Evidence missed at retrieval is unrecoverable at generation, and broadening the retrieved set introduces noise and attention dilution. Performance is further bounded by parser and embedding quality, while per-passage similarity scoring weakens when answers require aggregating weakly relevant pages or following multi-hop cross-references \cite{wang2024knowledge, wu2025molorag}, pressures that motivate replacing the fixed pipeline with adaptive, multi-step reasoning.

\subsection{Agentic Pipelines}
\label{sec:arch-agent}

Agentic pipelines replace the fixed retrieve-then-generate control flow with iterative evidence acquisition. Drawing on the broader agentic-LLM framing \cite{plaat2025agentic, wang2024survey}, we class as agentic any MP-VRDU system that performs multi-pass evidence acquisition or distributes work across specialised agents whose call sequence adapts at inference; multi-component systems with predetermined execution order, even when described as ``agents'', remain within the retriever-generator family. One or more controller MLLMs decide what to inspect next from intermediate observations and issue tool calls for retrieval, navigation, cropping, or OCR \cite{zheng2026docv, sun2025docagent}. Roles such as retrieval, inspection, synthesis, and adjudication may be distributed across agents, each backed by a different backbone selected for its subtask \cite{han2025mdocagent} rather than tied to a single generator. This design lifts the single-pass bottleneck, allows compute to scale with question difficulty rather than document length, and renders the trajectory itself an inspectable audit trail. The family is well suited to questions whose evidence is distributed across pages or modalities and to multi-hop questions that benefit from query reformulation between steps \cite{wang2025vidorag, sun2025docagent}, as well as visually heterogeneous documents whose text, tables, and figures benefit from tool-specific inspection \cite{han2025mdocagent, zhu2025doclens}. The flexibility carries costs. The variable-length action sequence makes inference cost difficult to bound before execution, controllers remain sensitive to prompt phrasing and may commit early to wrong evidence with silent failure, reproducibility is complicated by stochastic generation, and the pipeline inherits the recall ceiling of every retriever, parser, or OCR module it calls.

\section{Multi-Page Multimodal Representation}
\label{sec:modality}

While Section~\ref{sec:architecture} categorises MP-VRDU systems by where cross-page evidence is processed, this section analyses the problem from a modality perspective. Single-page VRDU already handles text, vision, and layout through the well-established mechanisms outlined in Section~\ref{sec:background}. The multi-page setting introduces a shared pressure: the full document often exceeds the backbone's context window, and even when it fits, the answer occupies only a small fraction of a long input where attention is diluted and key evidence is easily lost in the middle \cite{liu2024lost}. This forces each modality to trade page coverage against reading fidelity, with the resolutions differing across text, vision, and structure.

\subsection{Text Modality}
\label{sec:modality-text}

Text is the most information-dense modality in MP-VRDU but also the most heterogeneous in volume across pages and the most exposed to errors that compound when extracted across tens or hundreds of pages. Text reaches the backbone through two acquisition paths that recur across architectures. OCR-extracted strings support cheap indexing, natural answer-passage attribution, and easy bridging to visual elements through captions, but propagate parser errors and lose chart, layout, and handwriting semantics; pixel-recovered text read directly by an OCR-free MLLM avoids parser dependence but demands high-resolution input and large-scale text-recognition pretraining that scale poorly with document length \cite{kim2021donut,ye2023ureader}. Three patterns have emerged for handling text at document scale, each correlated with a different architectural family. \textbf{Per-page aggregation} encodes each page in detail and combines per-page representations into compact cross-page summaries through attention or memory tokens, typical of page-to-document encoders such as Hi-VT5 \cite{tito2023hierarchical}, GRAM \cite{blau2024gram}, and Arctic-TILT \cite{borchmann2025arctictilt}. \textbf{Whole-document compression} retains all pages within a single backbone context by aggressively reducing per-page text tokens through resamplers, layout-aware compressors, or token-reduction modules, typical of long-context backbone-centric MLLMs such as mPLUG-DocOwl2 \cite{hu2025mplugdocowl2} and TextHawk2 \cite{yu2024texthawk2}. \textbf{Retrieve-then-read} abandons full-document encoding entirely, indexing text at chunk or page level and passing only a retrieved subset to the reader, characteristic of retriever-generator and agentic systems such as CREAM \cite{zhang2024cream}, MultiDocFusion \cite{shin2025multidocfusion}, DocAgent \cite{sun2025docagent}, and MDocAgent \cite{han2025mdocagent}; agentic variants extend this into multiple selection-reading rounds. Each pattern carries a characteristic risk: per-page aggregation discards token-level detail at the summary step, whole-document compression sacrifices fine-grained reading on individual pages, and retrieve-then-read makes evidence missed at retrieval unrecoverable at the reading stage. The last two routes also depend on document-level structure to decide which pages or sections to read in detail, an issue we return to in Section~\ref{sec:modality-structure}.

\subsection{Visual Modality}
\label{sec:modality-visual}

Visual content in MP-VRDU spans text-bearing structures such as dense paragraphs and tables to non-textual evidence such as charts, figures, and handwriting that OCR cannot recover. Visual encoding faces a sharper version of the coverage-fidelity tension than text, since token cost scales jointly with page count and per-page resolution while the detail required to read charts, table cells, and handwriting cannot survive aggressive downsampling. Three patterns have emerged for managing this tension. \textbf{Aggressive token compression} keeps every page available in context but reduces per-page visual tokens through resamplers or convolutional reducers, as in mPLUG-DocOwl2 \cite{hu2025mplugdocowl2} and TextHawk2 \cite{yu2024texthawk2}; this preserves coverage at the cost of per-page detail, and is often paired with OCR text to recover fine reading where the downsampled image is insufficient. \textbf{Selective high-resolution reading} restricts expensive encoding to a small subset of pages surfaced by retrieval or controller-driven inspection, sending those high-fidelity tokens into the reader while leaving other pages unread, as in Doc-V* \cite{zheng2026docv} and SimpleDoc \cite{jain2025simpledoc}. \textbf{Retrieval-only visual embeddings} decouple encoding from reading altogether: page images are encoded purely to produce retrieval embeddings, and detailed reading of selected pages is handed to a separate text or vision module, as in M3DocRAG \cite{cho2024m3docrag} and VDocRAG \cite{tanaka2025vdocrag}. The patterns are not mutually exclusive in practice; several systems pair low-resolution thumbnails for navigation with selective high-resolution inspection of pages a retriever or tool call surfaces \cite{zheng2026docv,jain2025simpledoc}. Each pattern nonetheless carries a characteristic risk: aggressive compression loses fine text and small visual elements, selective high-resolution reading inherits the recall ceiling of whichever tool selects the pages, and retrieval-only visual embeddings may rank layout-similar pages without confirming that the answer-bearing region is present.

\subsection{Cross-Page Structure}
\label{sec:modality-structure}

Within-page mechanisms inherited from single-page VRDU, including 2D positional encodings, bounding-box coordinates, and reading-order signals attached to individual tokens, describe only per-page geometry and do not extend to relationships between sections, pages, or references. The structural problem distinctive to MP-VRDU is \textit{cross-page} structure: section hierarchies that span multiple pages, information referenced or continued across page boundaries, and figure-to-text or table-to-caption linkages that no longer fit within a single page's view. Two routes address it, with the second now dominant. Some backbone-centric systems train the model to predict cross-page structure end-to-end through multi-page parsing or document-level pretraining objectives \cite{hu2025mplugdocowl2,duan2025docopilot}, but pretraining supervision at document scale remains scarce and expensive to obtain, limiting how far this approach can reach in practice. The dominant route exposes parsed structural information to the model directly, and three patterns recur within it. \textbf{Structure-aware chunking} uses a parser to extract section trees and heading hierarchies, then segments documents along semantic boundaries so retrieved units respect section scope rather than sliding-window cuts, as in MultiDocFusion \cite{shin2025multidocfusion}. \textbf{Hierarchical retrieval indices} route queries first to coarse units such as sections or chapters and then to fine-grained chunks within the selected branch, reducing irrelevant pages while preserving cross-page context, as in MHier-RAG \cite{gong2025mhierrag} and MLDocRAG \cite{zhang2026mldocrag}. \textbf{Graph-based traversal} converts headings, entities, and cross-references into a knowledge graph, allowing controllers to follow typed edges across pages rather than relying on dense similarity alone, as in KGP \cite{wang2024knowledge}. All three depend on parser quality: errors in heading detection, table boundary segmentation, or reading-order extraction propagate directly into every downstream retrieval and navigation step, and reliable parsers themselves remain tuned for single-page or template-bound documents rather than the heterogeneous long-form documents on which MP-VRDU is evaluated.

\section{Inference-Time Adaptation}
\label{sec:inference-time}

Inference-time adaptation refers to strategies that improve MP-VRDU performance through changes to the inference pipeline rather than to model parameters. A frozen language or vision-language backbone is paired with off-the-shelf retrievers, structural parsers, prompted reasoning loops, or multi-agent orchestrations, none of which require gradient updates on the backbone itself. This approach is particularly suited to MP-VRDU settings where supervised data is scarce and expensive to collect, where backbones must remain exchangeable across systems, or where rapid adaptation to new document types is needed without retraining. We organise these strategies along the inference pipeline: evidence selection through retrieval and navigation (Section~\ref{sec:tf-retrieval}), reasoning over the assembled evidence (Section~\ref{sec:tf-reasoning}), and coordination across multiple frozen components (Section~\ref{sec:agent-orchestration}).

\subsection{Retrieval and Navigation}
\label{sec:tf-retrieval}

Inference-time retrieval selects a compact subset of pages, chunks, or graph nodes for the reader using off-the-shelf encoders, structural parsers, or LLM-based rerankers, without fine-tuning the retriever on document data. Retrieval quality bounds every downstream reasoning step, so the choice of signal dominates system behaviour on long documents. Two broad families have emerged, distinguished by whether passages are scored in isolation or navigated through explicit relations.

\paragraph{Similarity-based retrieval.} Similarity-based methods score each passage against the query independently, using sparse lexical signals, dense embeddings, or combinations of both. The dominant dense signal in MP-VRDU is vision-language late-interaction encoding such as ColPali \cite{faysse2024colpali} or its variants, which extends the late-interaction paradigm of ColBERT \cite{khattab2020colbert} to page images by producing token-level query vectors and patch-level page embeddings over the rendered page, then scoring a query-page pair by taking the maximum cosine similarity of each query token over all page patches and summing the result. This MaxSim aggregation preserves fine-grained match between specific query terms and specific page regions, which is critical when a single question token may need to align with one figure caption, table cell, or paragraph buried in a long document. M3DocRAG \cite{cho2024m3docrag} retrieves at chunk and page granularity with this signal alone, while AVIR \cite{li2025avir} adds an adaptive selector that examines the per-query score distribution and either thresholds at a fixed cutoff for short documents or applies K-means clustering with a top-$K$ bound on longer ones, scaling the number of returned pages with the score distribution rather than using a static $K$. Joint variants combine signals at the reranking stage: CREAM \cite{zhang2024cream} pairs dense retrieval with iterative LLM-based reranking across overlapping chunk groups, SimpleDoc \cite{jain2025simpledoc} combines ColPali embeddings with LLM-generated summary reranking under persistent working memory that refines the joint score across rounds, and MDocAgent \cite{han2025mdocagent} runs parallel text and image retrieval tracks whose outputs are aggregated by a downstream agent rather than fused at encoding time. Independent scoring handles paraphrase and cross-modal queries well, but cannot follow cross-page references or multi-hop dependencies that no single passage carries.

\paragraph{Structure-based retrieval.} Structure-based methods exploit a parsed representation of the document to make passage relations first-class objects. Tree-based variants such as MultiDocFusion \cite{shin2025multidocfusion} and DocAgent \cite{sun2025docagent} extract section hierarchies and outlines and group passages along semantic boundaries so retrieved chunks preserve parent context, while MHier-RAG \cite{gong2025mhierrag} extends this by routing queries through dual in-page and cross-page indices with LLM-based reranking across both. Graph-based variants treat headings, entities, and cross-references as graph nodes with typed edges, allowing controllers to traverse from one passage to related ones rather than relying on similarity alone. KGP \cite{wang2024knowledge} constructs a knowledge graph over passages and document-structure nodes such as pages, paragraphs, and tables with edges encoding lexical similarity and structural belonging, then uses an instruction-tuned LLM as a traverser that predicts the next evidence node from the visited path. MoLoRAG \cite{wu2025molorag} builds a page graph in which edge candidates are filtered by embedding similarity and then scored for logical relevance by a frozen vision-language prompt, combining semantic and logical signals for multi-hop traversal. MLDocRAG \cite{zhang2026mldocrag} constructs the graph per query, anchoring chunk nodes to LVLM-generated query nodes so the structure adapts to the question rather than being fixed at index time. Structure-based methods recover references that similarity scoring misses, at the cost of construction overhead, dependence on parser quality, and traversal cost that grows with hop count.

\subsection{Inference-Time Reasoning Strategies}
\label{sec:tf-reasoning}

Inference-time reasoning strategies govern how a frozen language or vision-language model accesses, updates, and reasons over evidence during decoding, without modifying model parameters. In MP-VRDU, these strategies can be broadly categorised by whether evidence acquisition remains static or becomes dynamically interactive during inference. \textbf{Closed-buffer reasoning} operates over a fixed evidence set assembled before generation begins, whereas \textbf{tool-interactive reasoning} allows the model to iteratively acquire additional evidence through retrieval or page-level actions during the reasoning process.

\paragraph{Closed-buffer reasoning.} Closed-buffer reasoning performs multi-step reasoning over a fixed evidence buffer assembled by a prior retrieval step or by a compressed document representation. Chain-of-thought prompting \cite{wei2022chainofthought} is the canonical instantiation, eliciting intermediate reasoning steps that integrate textual, tabular, and visual evidence before producing the final answer. MHier-RAG \cite{gong2025mhierrag} applies structured chain-of-thought prompts over hierarchical retrieval outputs, although a single reasoning chain remains vulnerable to compounding errors on dense layouts. SimpleDoc \cite{jain2025simpledoc} adds self-refinement and self-verification prompts that emit revised queries when the assembled evidence is judged insufficient, but the loop terminates on the backbone's own judgement and inherits its blind spots. Closed-buffer reasoning is fundamentally bounded by retrieval completeness: evidence absent from the buffer cannot be recovered by subsequent reasoning or verification, regardless of how many self-checks are applied.

\paragraph{Tool-interactive reasoning.} Tool-interactive reasoning interleaves reasoning and evidence acquisition during inference, allowing newly observed pages to guide subsequent reasoning steps. The ReAct paradigm \cite{yao2022react} is the canonical instantiation, alternating intermediate reasoning traces with tool calls whose outputs feed the next reasoning stage. Doc-React \cite{wu2025docreact} instantiates this with flat retrieval, InfoNCE-guided sub-query refinement, and entropy-based stopping, although each iteration still selects from a single retriever's ranked list and recovery from a poor encoder match relies on subsequent query rewrites. Outline-driven variants combine ReAct with the structured navigation of Section~\ref{sec:tf-retrieval}: DocAgent \cite{sun2025docagent} traverses a tree-formatted XML outline through section-level search and fetch tools chosen turn-by-turn by an actor agent, inheriting the parser's outline quality and degrading on documents without exploitable structure. A complementary refinement is coarse-to-fine visual reading, in which fetched pages are inspected at multiple resolutions rather than at a fixed scale so that high-resolution encoding is reserved for regions surfaced by the coarse pass, concentrating compute on task-relevant regions but making reasoning conditional on the localiser since the fine pass cannot recover when the coarse pass misroutes. Adaptive evidence acquisition improves localisation and supports more exploratory multi-page reasoning, but introduces tool latency that grows with rounds, sensitivity to prompt design and to early action choices that subsequent steps cannot recover from, and attention dilution within fetched pages at their reading resolution.

\subsection{Agentic Orchestration}
\label{sec:agent-orchestration}

Retriever-generator and agentic pipelines often distribute a single MP-VRDU question across multiple specialised language or vision-language model calls, each handling a particular modality, reasoning step, or verification role. The frozen backbones are not modified; coordination is achieved through prompt design and call ordering. The two families correspond directly to whether the trajectory of calls is fixed at design time or adapts to the question at inference: retriever-generator pipelines run a trajectory fixed in advance regardless of how many agents they invoke, while agentic pipelines treat the call sequence itself as a variable.

\paragraph{Fixed-trajectory pipelines.} Fixed-trajectory pipelines follow a control flow set in advance: which calls are made, in what order, and how their outputs are combined are all determined at design time. Modality-partitioned variants dispatch readers by input type and aggregate their outputs through a final synthesis step. MDocAgent \cite{han2025mdocagent} runs five parallel agents in which a general agent provides initial multimodal context, a critical agent extracts key evidence, dedicated text and image agents read their respective tracks, and a summariser fuses the four outputs into a single answer. M2RAG \cite{duan2025m2rag} splits work between a text filter, a visual extractor, and a modal fuser that applies a confidence-conditional rule, falling back to one modality when the other reports no information and fusing only when both contribute. Such pipelines improve coverage relative to single-agent baselines, but lack an error-correction stage: a wrong choice of routing or ordering cannot be revised mid-execution, and an incorrect output from one branch can dominate aggregation when the other branch lacks counter-evidence.

\paragraph{Adaptive-trajectory pipelines.} Adaptive-trajectory pipelines treat the call sequence itself as a variable: which agents are invoked and how many rounds the loop runs are decided per question rather than fixed in advance. Loop-based variants such as ViDoRAG \cite{wang2025vidorag} cycle a seeker that retrieves coarse candidates, an inspector that reviews them at fine resolution and emits reflective feedback, and an answerer that performs a consistency check, with a Gaussian mixture model setting per-query bounds from the similarity distribution rather than using a static top-$K$. Proposer-verifier variants insert an explicit verification stage between candidate generation and final answer: DocLens \cite{zhu2025doclens} samples multiple candidates and routes them through an adjudicator that compares reasoning paths and selects the most consistent conclusion rather than fusing them unconditionally, while DocAgent \cite{sun2025docagent} pairs an actor with a reviewer that cross-checks through complementary tools and writes disagreements into a shared reflection memory that updates the actor's behaviour on subsequent rounds. MACT \cite{yu2025visual} decomposes the loop further into planning, execution, judgement, and answer agents, separating the judge from the corrector to avoid the same agent excusing its own errors. Adaptive-trajectory pipelines improve the accuracy-latency trade-off, but inference cost remains hard to bound before execution and the verifier typically draws on the same backbone family as the proposer, so shared blind spots can persist when both agents fail in the same way.

\section{Training Strategies}
\label{sec:training}

While Section~\ref{sec:inference-time} held the backbone fixed, many MP-VRDU systems extend these pipelines by training specific components. We organise MP-VRDU training by the inference-time component being trained, namely the retriever, the reasoner, and the agentic orchestration policy. Two cross-cutting techniques, parameter-efficient tuning and multi-stage curriculum training, recur across all three targets.

\subsection{Trained Retrieval Components}
\label{sec:train-retrieval}

Off-the-shelf encoders such as ColPali \cite{faysse2024colpali} score pages in isolation, leaving three gaps that retriever training aims to address. Modality-alignment variants pretrain a visual document encoder against OCR text: VDocRAG \cite{tanaka2025vdocrag} introduces representation-compression retrieval and representation-compression generation objectives that compress page images into single-vector embeddings while preserving alignment with their OCR-derived textual content, retaining text alignment without the storage and latency cost of late-interaction patch embeddings. Cross-page-fusion variants freeze the encoder and learn a neighbour-blending head: DFVC \cite{qian2026accurate} adds a multi-layer perceptron adapter with gating and residual connections that fuses each page embedding with its neighbours, so evidence spread across adjacent pages, which independent per-page scoring would split, becomes retrievable as a unit. Logical-relevance variants train a small scorer over a page graph: MoLoRAG \cite{wu2025molorag} fine-tunes a frozen Qwen2.5-VL scorer to rank neighbour pages during graph traversal, supplying the cross-page logical signal that embedding similarity alone often misses and effectively turning retrieval into a learned multi-hop walk rather than a single-step similarity match.

\subsection{Trained Reasoning Components}
\label{sec:train-reasoning}

Prompted reasoning leaves quality sensitive to prompt design and base-model behaviour, motivating approaches that train the reasoning component directly. Reward-driven variants train a reasoning policy with explicit evidence supervision: DocR1 \cite{xiong2026docr1} fine-tunes a vision-language backbone with Evidence-Page-Guided GRPO, whose composite reward combines format compliance, answer accuracy, and a page-citation indicator that explicitly credits trajectories which cite the correct evidence pages, producing structured outputs that interleave thinking traces, evidence pages, and answers. Trace-imitation variants instil a reasoning chain through supervised fine-tuning and preference optimisation: Chain-of-Reading \cite{guo2026endtoend} bakes in a four-stage plan-execute-integrate-verify chain through a staged pipeline of LoRA supervised tuning on a mixed corpus, full-parameter tuning on a curated reasoning dataset, and direct preference optimisation on preference pairs that prefer faithful chains over hallucinated shortcuts. Uncertainty-driven variants train an aggregation head rather than a trajectory: DocSLM \cite{hannan2025docslm} learns an entropy-based calibrator that selects the lowest-uncertainty answer across streaming page segments, providing a self-consistency signal that scales to long documents under tight parameter budgets without requiring full-document context.

\subsection{Trained Agentic Orchestration}
\label{sec:train-orchestration}

Prompt-only orchestration is unreliable on smaller open-weight backbones that cannot follow ReAct protocols. Behavioural-cloning variants distil agent trajectories from a stronger model: DocDancer \cite{zhang2026docdancer} fine-tunes a Qwen3 agent on synthesised exploration-then-synthesis traces that pair \texttt{Search} and \texttt{Read} ReAct actions with intent-conditioned tool selection, training the smaller model to match the trajectory shape of a larger teacher rather than learning the underlying policy from scratch. Reinforcement-learned variants reward correct tool use and trajectory shape: Doc-V* \cite{zheng2026docv} couples supervised fine-tuning on coarse-to-fine ReAct traces with GRPO over a two-action space of \texttt{retrieval\_page} and \texttt{fetch\_page} calls plus entropy-based stopping that learns when to commit to an answer rather than fetch another page, while MACT \cite{yu2025visual} trains a multi-agent policy with agent-wise adaptive test-time scaling that allocates sample budgets across the planner, executor, judge, and answerer roles per question, spending more samples on the role that the question's difficulty most stresses rather than running every agent at a fixed budget.

\subsection{Cross-Cutting Training Techniques}
\label{sec:train-crosscutting}

Two techniques recur across retriever, reasoner, and orchestrator training. Full backbone fine-tuning is prohibitive at long-document context lengths and erases pretrained capabilities, so the dominant pattern freezes the backbone and trains lightweight modules through parameter-efficient tuning approaches such as LoRA \cite{hu2022lora}, projection heads, or task-specific adapters. LayTokenLLM \cite{zhu2025simple} freezes the language model and adds a layout tokenisation module plus LoRA, encoding spatial layout through compressed positional tokens with negligible context overhead. CREAM \cite{zhang2024cream} applies the LLaMA-Adapter-V2 recipe above a frozen Pix2Struct encoder, training only LoRA matrices, attention pooling, and bias and normalisation parameters in its answerer. The same pattern appears in DFVC's frozen-ColPali adapter and DocDancer's LoRA trajectory tuning, confirming that the choice between parameter-efficient and full fine-tuning is more often a compute-budget decision than a target-component one. Separately, several systems converge on multi-stage curricula that start on single-page or local objectives before introducing document-scale supervision, exploiting the fact that text reading and layout understanding transfer from single-page data while cross-page integration must be learned separately. mPLUG-DocOwl2 \cite{hu2025mplugdocowl2} uses three stages of single-image structure learning on DocStruct4M, MP-DocStruct1M continued pretraining, and multi-task fine-tuning; DocSLM \cite{hannan2025docslm} extends this to a five-stage schedule that adds image-OCR alignment at the start and streaming-abstention calibration at the end, demonstrating that the curriculum can be lengthened to cover both pre-compression alignment and post-reasoning uncertainty calibration without abandoning the staged structure.

\section{Datasets}
\label{sec:datasets}

Datasets in MP-VRDU span three roles: pretraining, fine-tuning, and benchmarking. Pretraining historically relied on single-page archives such as IIT-CDIP \cite{lewis2006building} and OCR-IDL \cite{biten2022ocridl}, with genuinely multi-page corpora only recently emerging through DocStruct4M's mixed single- and multi-page pretraining \cite{hu2024mplugdocowl} and the purely multi-page MP-DocStruct1M, Doc-750K, and scientific PaperPDF \cite{hu2025mplugdocowl2,duan2025docopilot,xie2024wukong}. Fine-tuning has shifted toward bespoke synthetic generation per model, with strong vision-language models writing question-answer or trajectory data over scraped PDFs since per-page annotation does not scale to long documents \cite{tanaka2024instructdoc,jain2025simpledoc,xiong2026docr1,zhang2026docdancer}. Benchmarks have evolved from mid-length extractive evaluation on MP-DocVQA, DUDE, and SlideVQA \cite{tito2023hierarchical,vanlandeghem2023document,tanaka2023slidevqa} to long-context multimodal evaluation on MMLongBench-Doc, LongDocURL, and DocBench \cite{ma2024mmlongbenchdoc,deng2025longdocurl,zou2024docbench}, with retrieval and page-localisation increasingly scored alongside answer accuracy \cite{hui2024uda,xiong2026docr1,zhu2025doclens}. Five benchmarks (MP-DocVQA, MMLongBench-Doc, DUDE, DocVQA, and InfographicVQA \cite{mathew2022infographicvqa}) appear in roughly half or more of recent MP-VRDU papers and constitute the de-facto axes for cross-model comparison. A detailed catalogue with per-dataset statistics is provided in Table~\ref{tab:dataset_survey}.

\section{Open Gaps}
\label{sec:gaps}

The preceding sections show that MP-VRDU has produced a rich landscape of architectures, modality-handling strategies, inference-time compositions, and training recipes. Despite this rapid progress, three persistent gaps cut across the four architectural families and remain unresolved by any current approach.

\subsection{Data Scarcity and Evaluation Integrity}

Per-page annotation does not scale to documents averaging eight to twenty pages, so most systems rely on vision-language-model-synthesised question-answer pairs and trajectories generated over scraped PDFs \cite{tanaka2024instructdoc,jain2025simpledoc,xiong2026docr1,zhang2026docdancer}. Of the twelve papers in this review that introduce a synthesised corpus, nine source their PDFs from existing evaluation benchmarks, and MP-DocVQA alone appears in 15 fine-tuning mixes and 17 evaluation suites among the 36 surveyed papers, with DUDE and DocVQA showing similar overlap. The result is a contamination loop in which a model trains on a strong vision-language model's behaviour over benchmark PDFs and is then evaluated against those same PDFs under metrics often judged by the same model family, so performance reported as zero-shot routinely includes the test corpus in its training history. No paper in the surveyed corpus reports synthetic-data quality audits, document-hash de-duplication against test pools, or backbone-level contamination checks, while agentic stochasticity and per-paper benchmark curations further block direct cross-system comparison. Dominant benchmarks also remain overwhelmingly English-centric, leaving non-English layouts, scripts, and document conventions underrepresented in both training and evaluation. These limitations leave open the design of held-out closed-test benchmarks, mandatory contamination declarations covering pretraining, fine-tuning, and backbone leakage, generator-evaluator separation when synthetic data is involved, trajectory-level metrics paired with retrieval and grounding scores, and multilingual corpora that cover non-English layouts and conventions.

\subsection{The Evidence Pipeline Bottleneck}

MP-VRDU systems must reconcile high-resolution per-page reading with document-scale evidence aggregation, an objective that exceeds any MLLM's compute and context budget \cite{ma2024mmlongbenchdoc}. Every architectural family takes a lossy compromise, but the compromise sits at a different stage of the pipeline. Page-to-document encoders lose fine-grained evidence at the per-page compression step, since summary or memory tokens cannot retain every token-level detail. Backbone-centric MLLMs dilute attention across long contexts and lose evidence buried in the middle of the input \cite{liu2024lost}. Retriever-generator pipelines drop everything outside the top-$K$ and depend on the assumption that relevance is rankable by query similarity in a single step. Agentic systems inherit the recall ceiling of every tool they call, so weaknesses in any retriever or parser bound the entire trajectory regardless of how many rounds the controller runs. Compression is also typically fixed at index time and independent of the question, so evidence is discarded before the question is seen, and retrieval-stage chunking errors, OCR artefacts, parser failures, and per-passage scoring that ignores cross-references propagate directly into the answer. Iterative retrieval lifts this ceiling partially yet inherits the recall of every tool it calls. The bottleneck is fundamental: across page-to-document, backbone-centric, retriever-generator, and agentic families, no architecture has decoupled answer quality from a coverage-versus-fidelity trade-off, only shifted where the trade-off is made. Query-conditional and modality-aware compression that retains evidence based on the question's evidential demands, and retrieval that follows cross-references and logical relations rather than scoring passages in isolation, remain open directions.

\subsection{Cross-Page Structural Reasoning}

Within-page layout is largely inherited from single-page VRDU through 2D positional embeddings and spatial pretraining (Section~\ref{sec:background}), while cross-page structure such as section hierarchies that span multiple pages, page order, multi-page tables, footnote-to-text linkages, and figure-caption-text references spread across pages remains the distinctive multi-page problem. Three current routes each carry a distinct cost. Layout-as-objective approaches train the backbone to predict cross-page structure end-to-end through multi-page parsing or document-level pretraining objectives \cite{hu2025mplugdocowl2,duan2025docopilot}, but demand supervision that is expensive at document scale, since reliable cross-page structural annotation is difficult to collect at the volumes pretraining requires. Implicit pixel-based recovery depends on backbone capability and high-resolution input that aggressive downsampling readily compromises, so the same compression pressures discussed in Section~\ref{sec:modality-visual} weaken the structural signal whenever pages are compressed to fit document-scale inputs. Explicit parsed structure inherits parser fragility on headings, table boundaries, and reading order, with errors propagating into every downstream retrieval and navigation step; the parsers themselves remain tuned for single-page or template-bound documents rather than the heterogeneous long-form documents on which MP-VRDU is evaluated, and rarely recover relationships such as theorem-proof or footnote-reference linkages that cross page boundaries. No current route addresses the structural problem end-to-end without falling back on parsers tuned for single-page or template-bound documents. Cheaper structural supervision, learned cross-page representations that survive compression, and end-to-end structure-aware retrieval that does not depend on offline parser output remain open directions.
\section{LLM Usage and Other Notes}

\paragraph{LLM Usage.}
Large language models, specifically Claude Opus 4.7 and GPT 5.5, were used throughout the preparation of this review for spelling and grammar correction, sentence-level rewording, reading-flow refinement, and other writing-quality tasks. All conceptual content, argumentation, and selection of cited works are my own. The LLM was used as a writing aid rather than as a source of claims, and all factual statements, citations, and quantitative figures have been verified against primary sources before inclusion. Intellectual responsibility for the final document rests with the author.

\paragraph{Relationship to a Companion Survey Paper.}
Sections~\ref{sec:architecture}--\ref{sec:challenges} were adapted from a survey paper on the MP-VRDU frontier of which I am the first author. The corresponding material was largely written by me prior to any co-author contributions to the survey. The adaptation in this review trims, re-tones, and reorients the survey material toward the methodology research direction motivated in \S\S\ref{sec:gaps}--\ref{sec:targets}; the remaining sections (Introduction, Background, Open Problems, Targets and Loci, and Conclusion) are original to this document.

\paragraph{Supporting Materials.}
The following supporting artefacts were produced during the preparation of this review:
\begin{itemize}
  \item Paper inventory and cross-reference spreadsheet (surveyed systems, dataset and benchmark coverage, evaluation metrics): \\ \url{https://docs.google.com/spreadsheets/d/1ZoPeIvnZKkHOQN8G3Fjj9gY8jnLImZvcYQrL36njS78/edit?usp=sharing}
  \item Project repository (literature review \LaTeX{} source, build script, bibliography, and supporting code): \\ \url{https://github.com/LeweiXu/MP-VRDU/}
  \item Overleaf project for the companion survey paper: \emph{URL pending}.
\end{itemize}

% Bibliography style is set inside acl.sty as acl_natbib.bst; we have replaced
% that file with IEEEtran's content so numeric IEEE-style entries are produced.
\bibliography{custom}

\appendix

\end{document}
